% Part: turing-machines
% Chapter: undecidability
% Section: unsolvability-decision-problem

\documentclass[../../../include/open-logic-section]{subfiles}

\begin{document}

\olfileid{tur}{und}{uns}
\olsection{সিদ্ধান্ত সমস্যাটি অসমাধানযোগ্য}

\begin{thm}
\ollabel{thm:decision-prob}
সিদ্ধান্ত সমস্যাটি অসমাধানযোগ্য: এমন কোনো টুরিং যন্ত্র~$D$ নেই যা
প্রথম-ক্রমের যুক্তিবিদ্যার কোনো !!a{sentence}~$!B$ নিবেশ হিসেবে থাকা
ফিতায় শুরু করলে $D$ শেষ পর্যন্ত থামে এবং $!B$ সিদ্ধ হলে ও কেবল তখনই~$1$,
আর অন্যথায়~$0$ নির্গত করে।
\end{thm}

\begin{proof}
ধরা যাক, সিদ্ধান্ত সমস্যাটি সমাধানযোগ্য; অর্থাৎ এমন একটি টুরিং যন্ত্র~$D$
আছে। তাহলে নিচের উপায়ে থামার সমস্যা সমাধান করা যেত। এমন একটি টুরিং
যন্ত্র~$E$ নির্মাণ করি যা টুরিং যন্ত্র~$M_e$-এর সংখ্যা~$e$ এবং নিবেশ~$w$
পেলে সংশ্লিষ্ট !!{sentence}~$!T(M_e,w)\lif !E(M_e,w)$ গণনা করে এবং
ফিতার বাঁদিকের প্রথম ঘর পড়তে পড়তে থামে। তখন যন্ত্র $E\frown D$ নিবেশ
$e$ ও~$w$ পেলে প্রথমে $!T(M_e,w)\lif !E(M_e,w)$ গণনা করবে এবং তারপর সেই
নিবেশে সিদ্ধান্ত-সমস্যার যন্ত্র~$D$ চালাবে। $D$ ঠিক তখনই নির্গম~$1$-সহ
থামে যখন $!T(M_e,w)\lif !E(M_e,w)$ সিদ্ধ; অন্যথায় নির্গম~$0$ দেয়।
\olref[ver]{lem:halt-if-valid} ও \olref[ver]{lem:valid-if-halt} অনুসারে
$!T(M_e,w)\lif !E(M_e,w)$ সিদ্ধ যদি এবং কেবল যদি $M_e$ নিবেশ~$w$-তে
থামে। ফলে $E\frown D$ নিবেশ $e$ ও~$w$ পেয়ে ঠিক তখনই নির্গম~$1$-সহ
থামত যখন $M_e$ নিবেশ~$w$-তে থামে; অন্যথায় নির্গম~$0$-সহ থামত। অর্থাৎ
$E\frown D$ থামার সমস্যা সমাধান করত। কিন্তু
\olref[hal]{thm:halting-problem} থেকে জানি, এমন কোনো টুরিং যন্ত্র থাকতে
পারে না।
% BN-SRC-188 TARGET-CORRECTION-BEGIN
\par\smallskip\noindent\textnormal{\small\textbf{উৎস-সংশোধন (BN-SRC-188):} হিমায়িত প্রমাণে তিনবার $E\concat D$ লেখা হয়েছে, যদিও আগের যন্ত্র-সংযোজন সংজ্ঞা অনুসারে ধারাবাহিকভাবে যুক্ত যন্ত্রের প্রতীক $E\frown D$; $\concat$ স্ট্রিং-সংযোজনের জন্য সংরক্ষিত। বাংলা পাঠে তিনটি ঘটনে $E\frown D$ রাখা হয়েছে।} % BN-SRC-188 TARGET-CORRECTION-END
\end{proof}

\begin{cor}\ollabel{cor:undecidable-sat}%
প্রথম-ক্রমের যুক্তিবিদ্যার কোনো নির্বিচার !!{sentence} সন্তোষণীয় কি না তা
অনির্ণেয়।
\end{cor}

\begin{proof}
  ধরা যাক, সন্তোষণীয়তা টুরিং যন্ত্র~$S$ দিয়ে নির্ণেয়। তাহলে সিদ্ধান্ত
  সমস্যাটি নিচের মতো সমাধান করা যেত। নিবেশ হিসেবে কোনো
  !!a{sentence}~$!B$ পেলে $!B$-কে এক ঘর ডানে সরাও। ঘর~$1$-এ ফিরে
  $\lnot$ প্রতীকটি লেখো।
% BN-SRC-189 TARGET-CORRECTION-BEGIN
\par\smallskip\noindent\textnormal{\small\textbf{উৎস-সংশোধন (BN-SRC-189):} হিমায়িত প্রমাণের প্রথম ঘটনে বাক্যনামটি $B$ লেখা হলেও একই বাক্যের পরের ঘটনে ও পুরো যুক্তিতে $!B$ ব্যবহৃত হয়েছে। বাংলা পাঠে অনুপস্থিত অতিভাষিক চিহ্নটি পুনঃস্থাপন করা হয়েছে।} % BN-SRC-189 TARGET-CORRECTION-END

  এবার টুরিং যন্ত্র~$S$ চালাও। সেটি শেষ পর্যন্ত ফিতায়~$1$ (যদি
  $\lnot !B$ সন্তোষণীয় হয়) অথবা~$0$ (যদি $\lnot !B$ অসন্তোষণীয় হয়)
  নির্গত করে থামে। ঘর~$1$-এ $\TMstroke$ থাকলে সেটি মুছে ফেলো; ঘর~$1$
  খালি থাকলে একটি~$\TMstroke$ লেখো; তারপর থামো।

  এই টুরিং যন্ত্রটি সবসময় থামে, এবং ঠিক তখনই এর নির্গম~$1$ যখন
  $\lnot !B$ অসন্তোষণীয়; অন্যথায় নির্গম~$0$। যেহেতু $!B$ সিদ্ধ যদি
  এবং কেবল যদি $\lnot !B$ অসন্তোষণীয়, তাই যন্ত্রটি ঠিক তখনই~$1$
  নির্গত করে যখন $!B$ সিদ্ধ, আর অন্যথায়~$0$। অর্থাৎ যন্ত্রটি সিদ্ধান্ত
  সমস্যা সমাধান করত।
\end{proof}

\begin{explain}
অতএব “$!B$ কি প্রথম-ক্রমের যুক্তিবিদ্যার সিদ্ধ !!{sentence}?”—এই প্রশ্নে
সবসময় সঠিক “হ্যাঁ” বা “না” উত্তর দেয় এমন কোনো টুরিং যন্ত্র নেই। তবে
এমন একটি টুরিং যন্ত্র \emph{আছে} যা সবসময় সঠিক “হ্যাঁ” উত্তর দেয়, কিন্তু
উত্তর “না” হলে আর থামে না। প্রথম-ক্রমের যুক্তিবিদ্যার যথার্থতা ও
সম্পূর্ণতা উপপাদ্য এবং !!{derivation}s কার্যকরভাবে গণনা করা যায়—এই
দুটি তথ্য থেকে তা পাওয়া যায়।
\end{explain}

\begin{thm}
  \ollabel{thm:valid-ce}%
  প্রথম-ক্রমের !!{sentence}s-এর সিদ্ধতা অর্ধ-নির্ণেয়: এমন একটি টুরিং
  যন্ত্র~$E$ আছে যা প্রথম-ক্রমের যুক্তিবিদ্যার কোনো !!a{sentence}~$!B$
  নিবেশ হিসেবে থাকা ফিতায় শুরু করলে $E$ ঠিক তখনই শেষ পর্যন্ত থেমে~$1$
  নির্গত করে যখন $!B$ সিদ্ধ, আর অন্যথায় থামে না।
\end{thm}

\begin{proof}
  প্রথম-ক্রমের যুক্তিবিদ্যার সম্ভাব্য সব !!{derivation}s একটি কার্যকর
  অ্যালগরিদমে একটির পর একটি উৎপন্ন করা যায়। যন্ত্র~$E$ তা-ই করে; এবং
  $\Proves !B$ দেখানো কোনো !!a{derivation} পেলে নির্গম~$1$-সহ থামে।
  যথার্থতা উপপাদ্য অনুসারে $E$ নির্গম~$1$-সহ থামলে তার কারণ
  $\Entails !B$। সম্পূর্ণতা উপপাদ্য অনুসারে $\Entails !B$ হলে
  $\Proves !B$ দেখানো একটি !!a{derivation} আছে। $E$ যেহেতু সম্ভাব্য সব
  !!{derivation}s নিয়মমাফিক উৎপন্ন করে, তাই শেষ পর্যন্ত
  $\Proves !B$ দেখানো একটি নিগমন খুঁজে পাবে এবং নির্গম~$1$-সহ থামবে।
\end{proof}

\end{document}
